\documentclass[../../../include/open-logic-section]{subfiles}

\begin{document}

\olfileid{sth}{ordinals}{ordtype}
\olsection{ترتیب دیاں قسماں دے طور اُتے ترتیبی عدد}

بدل رکھن دے اصول نال، تے ہُن $\ZFminus$ وچ کم کردیاں، اسیں آخر
اوہ نتیجہ ثابت کر سکدے آں جیہڑا ساڈا مقصد سی:

\begin{thm}\ollabel{thmOrdinalRepresentation}
ہر حُسنِ ترتیب والی ساخت اکّو اک ترتیبی عدد نال ہم ساخت اے۔
\end{thm}

\begin{proof}
منّ لو $\tuple{A, <}$ حُسنِ ترتیب والی ساخت اے۔
\olref[basic]{ordisoidentity} نال ایہ ودھ توں ودھ اک ترتیبی عدد نال
ہم ساخت اے۔ سو تضاد کڈھن لئی منّ لو $\tuple{A, <}$ \emph{کسے وی}
ترتیبی عدد نال ہم ساخت نہیں۔ اسیں پہلاں ’’$\tuple{A, <}$ نوں جِنّا ہو
سکے چھوٹا کراں گے‘‘۔ تفصیل ایہ اے۔ جے کوئی حقیقی مُڈھلا ٹکڑا
$\tuple{A_a, <_a}$ کسے وی ترتیبی عدد نال ہم ساخت نہیں، تاں ایس خاصیت
والا سبھ توں چھوٹا $a \in A$ اے۔ فیر منّ لو $B = A_a$ تے
$\mathord{\lessdot} = \mathord{<_a}$۔ نہیں تاں منّ لو $B = A$ تے
$\mathord{\lessdot} = \mathord{<}$۔

تعریف مطابق $B$ دا ہر حقیقی مُڈھلا ٹکڑا کسے ترتیبی عدد نال ہم ساخت اے،
تے اُتلی دلیل نال ایہ عدد اکّو اک اے۔ سو بدل رکھن دے اصول نال ہِیٹھلا
سیٹ موجود اے تے اک فنکشن اے:
\[
	f = \Setabs{\tuple{\beta, b}}{b \in B\text{ تے }\ordeq{\beta}{\tuple{B_b, \lessdot_b}}}
\]
تضاد پورا کرن لئی اسیں وکھاواں گے کہ $f$، $\alpha \to B$ دا اک
ہم ساخت نقشہ اے، کسے ترتیبی عدد $\alpha$ لئی۔

صاف اے کہ $\ran{f} = B$۔ تے \olref[iso]{lemordsegments} نال
$f$ ترتیب قائم رکھدا اے، یعنی $\gamma \in \beta$ اُدوں تے صرف
اُدوں جدوں $f(\gamma) \lessdot f(\beta)$۔ ایہ وکھاؤن لئی کہ $\dom{f}$
ترتیبی عدد اے، \olref[basic]{corordtransitiveord} نال اینا وکھاؤنا
کافی اے کہ $\dom{f}$ متعدی اے۔ سو $\beta \in \dom{f}$ مقرر کرو، یعنی
$\ordeq{\beta}{\tuple{B_b, \lessdot_b}}$ کسے $b$ لئی۔ جے $\gamma \in \beta$،
تاں $\gamma \in \dom{f}$، \olref[iso]{wellordinitialsegment} نال۔
عام صورت وچ $\beta \subseteq \dom{f}$۔
\end{proof}

ایہ نتیجہ سانوں ہِیٹھلی تعریف دین دی اجازت دیندا اے، جیہڑی اسیں
\olref[vn]{sec} توں دینی چاہندے ساں:

\begin{defn}
جے $\tuple{A, <}$ حُسنِ ترتیب والی ساخت اے، تاں اوہدی ترتیب دی قسم
$\ordtype{A, <}$ اوہ اکّو اک ترتیبی عدد $\alpha$ اے جیہدے لئی
$\ordeq{\tuple{A, < }}{\alpha}$۔
\end{defn}

نالے ایس تعریف توں دو چنگے اصول ملدے نیں:

\begin{cor}\ollabel{ordtypesworklikeyouwant}
جے $\tuple{A, <}$ تے $\tuple{B, \lessdot}$ حُسنِ ترتیب والیاں ساختاں نیں:
\begin{align*}
	\ordtype{A, <} = \ordtype{B, \lessdot}&\text{ اُدوں تے صرف اُدوں جدوں }\ordeq{\tuple{A, <}}{\tuple{B, \lessdot}}\\
	\ordtype{A, <} \in \ordtype{B, \lessdot}&\text{ اُدوں تے صرف اُدوں جدوں }\ordeq{\tuple{A, <}}{\tuple{B_b, \lessdot_b}}\text{ کسے ایہو جیہے }b \in B
\end{align*}
\end{cor}

\begin{proof}
برابری \olref[basic]{ordisoidentity} نال قائم اے۔ دوجی گل ثابت کرن لئی
منّ لو $\ordtype{A, <} = \alpha$ تے $\ordtype{B, \lessdot} = \beta$،
تے منّ لو $f \colon \beta \to \tuple {B, \lessdot}$ ساڈا ہم ساخت
نقشہ اے۔ فیر:
\begin{align*}
	\alpha \in \beta&\text{ اُدوں تے صرف اُدوں جدوں }\funrestrictionto{f}{\alpha} \colon \alpha \to B_{f(\alpha)}\text{ اک ہم ساخت نقشہ اے}\\
	&\text{ اُدوں تے صرف اُدوں جدوں }\ordeq{\tuple{A, <}}{\tuple{B_{f(\alpha)}, \lessdot_{f(\alpha)}}}\\
	&\text{ اُدوں تے صرف اُدوں جدوں }\ordeq{\tuple{A, <}}{\tuple{B_b, \lessdot_b}}\text{ کسے $b \in B$ لئی}
\end{align*}
\olref[iso]{ordisounique}، \olref[iso]{wellordinitialsegment} تے
\olref[basic]{ordissetofsmallerord} نال۔
\end{proof}

\end{document}
